% ================================================================
\section{The Arakelov Bridge: Algebraic Geometry Meets Spectral Theory}
\label{sec:arakelov}
% ================================================================

The Arakelov subsystem (\lean{Arakelov/*.lean} + \lean{Zeta/ArakelovBridge.lean},
5~files, $\sim$1{,}660~lines, 0~sorry, 1~axiom) establishes an algebraic-geometric
interpretation of the Gram matrix, connecting the Nyman-Beurling spectral
approach to Arakelov intersection theory on~$\Spec(\ZZ)$.

The central insight: the Gram matrix $G_N$ with entries
$G(j,k) = \int_0^1 \{1/(jx)\}\{1/(kx)\}\,dx$ decomposes as
\begin{equation}
  G(j,k) = \underbrace{\frac{\gcd(j,k)^2}{12jk}}_{\text{finite part } G_{\mathrm{fin}}}
  + \underbrace{\text{perturbation}}_{\text{archimedean part } G_{\mathrm{arch}}}
  \label{eq:arakelov-decomp}
\end{equation}
where the finite part is determined by the prime factorizations
of~$j$ and~$k$ (through the GCD), and the archimedean part captures
the analytic correction from higher Bernoulli terms.

\subsection{The Four Layers}

The Arakelov bridge is built in four layers, each with zero sorry:

\begin{center}
\small
\begin{tabular}{@{}cllrl@{}}
\toprule
\textbf{Layer} & \textbf{File} & \textbf{Content} & \textbf{Lines} & \textbf{Key Result} \\
\midrule
1 & \lean{WeilDivisor.lean}       & Divisors on $\Spec(\ZZ)$       & 219 & $\log\deg(D_k) = \log k$ \\
2 & \lean{ArithmeticDivisor.lean} & Arakelov metric                & 336 & $\deg(\hat{D}_1) = 0$ \\
3 & \lean{GramBridge.lean}        & Tropical $\to$ Gram connection & 366 & $\text{b1Entry} = \gcd^2/(12jk)$ \\
4 & \lean{ArakelovFusion.lean}    & Fusion theorem                 & 370 & $G = G_{\mathrm{fin}} + G_{\mathrm{arch}}$ \\
\bottomrule
\end{tabular}
\end{center}

\subsubsection*{Layer 1: Weil Divisors (\lean{WeilDivisor.lean})}

A Weil divisor on $\Spec(\ZZ)$ is a formal sum
$D = \sum_p v_p(D) \cdot [p]$ of prime ideals with integer multiplicities.
The \emph{log-degree} is $\log\deg(D) = \sum_p v_p(D) \cdot \log p$.
For the factorization divisor of an integer~$k$:
$\log\deg(D_k) = \log k$
(\lean{logDegree\_factorizationDivisor}).

\subsubsection*{Layer 2: Arithmetic Divisors (\lean{ArithmeticDivisor.lean})}

An arithmetic divisor $\hat{D} = (D, g)$ pairs a Weil divisor with a
Green's function $g \in \RR$ (the archimedean contribution). The
\emph{arithmetic degree} is $\deg(\hat{D}) = \log\deg(D) + g$.
For the Nyman-Beurling divisor encoding with $g_k = -\log k$:
\[
  \deg(\hat{D}_k) = \log k + (-\log k) = 0
\]
(\lean{bdDivisor\_one\_degree\_zero}). This is the
\emph{product formula} $|k|_{\mathrm{fin}} \cdot |k|_\infty = 1$
expressed as arithmetic degree zero.

\begin{correspondence}[Hodge Index Theorem]
In Arakelov geometry, the Hodge Index Theorem states that the
arithmetic intersection pairing is negative definite on divisors
of degree zero modulo numerical equivalence. Our product formula
shows the BD divisors have degree zero, placing them in the regime
where Hodge Index applies. The Gram matrix's positive-definiteness
failure (i.e., $\lambda_{\min}(G_N) \to 0$) is precisely the
Hodge Index Theorem manifesting in the spectral domain.
\end{correspondence}

\subsubsection*{Layer 3: The Gram Bridge (\lean{GramBridge.lean})}

The tropical intersection (\lean{gcdIntersection}) computes
\[
  \langle D_j, D_k \rangle_{\mathrm{trop}} = \sum_p \min\bigl(v_p(j), v_p(k)\bigr) \cdot \log p
  = \log\gcd(j,k)
\]
The B$_1$ Arakelov entry is then
$\text{b1ArakelovEntry}(j,k) = \gcd(j,k)^2 / (12jk)$
(\lean{b1ArakelovEntry}), which equals the exponential of
twice the tropical intersection:
\[
  G_{\mathrm{fin}}(j,k) = \frac{e^{2\langle D_j, D_k \rangle_{\mathrm{trop}}}}{12 \cdot j \cdot k}
\]
(\lean{b1\_from\_tropical}).

\subsubsection*{Layer 4: The Fusion (\lean{ArakelovFusion.lean})}

The Fusion theorem (\lean{arakelov\_eq\_skeleton}) proves that the
Arakelov road (\lean{GramBridge}) and the Cathedral road
(\lean{BernoulliSkeleton}) independently computed the \emph{same}
function: $\text{b1ArakelovEntry} = \text{b1Entry}$ (definitional equality).
This enables:
\begin{enumerate}
  \item \textbf{PSD Transfer} (\lean{arakelov\_pairing\_psd}):
    Smith's 1876 theorem (GCD matrices are PSD) transfers to the
    Arakelov pairing via the definitional bridge. Effective arithmetic
    divisors have non-negative self-intersection.

  \item \textbf{Gram Decomposition} (\lean{gram\_arakelov\_decomposition}):
    $G(j,k) = G_{\mathrm{fin}}(j,k) + G_{\mathrm{arch}}(j,k)$,
    with concrete witnesses for both parts.

  \item \textbf{GCD Structure} (\lean{G\_fin\_gcd\_structure}):
    the finite part depends only on $\gcd(j,k)$ and the product~$jk$.
    On the diagonal, $G_{\mathrm{fin}}(j,j) = 1/12$ for all~$j$
    (\lean{G\_fin\_diagonal}).
\end{enumerate}

\subsection{The Tropical-Spectral Pipeline}

The Arakelov bridge establishes a complete pipeline:
\[
  \text{Tropical geometry}
  \xrightarrow{\text{GCD}}
  \text{Arakelov intersection}
  \xrightarrow{\exp}
  \text{Gram matrix}
  \xrightarrow{\lambda_{\min}}
  \text{RH}
\]

\begin{principle}[Algebraic Geometry Dictates Spectral Decay]
The rate at which $\lambda_{\min}(G_N) \to 0$ is controlled by
the Arakelov intersection pairing of the BD divisors. The finite
part $G_{\mathrm{fin}}$ (GCD structure) provides the skeleton---a
PSD matrix that Smith's 1876 theorem protects from below. The
archimedean part $G_{\mathrm{arch}}$ (analytic correction) provides
the refinement that drives the eigenvalue toward zero at the rate
demanded by the Riemann Hypothesis.
\end{principle}

\subsection{Remaining Gap}

The sole remaining axiom on the Arakelov road is
\lean{hodge\_index\_eigenvalue\_bound} (in \lean{ArakelovBridge.lean}),
asserting $\lambda_{\min}(G_N) \leq C/N^{1+\varepsilon}$. Graduating
this axiom requires:
\begin{enumerate}
  \item The M\"obius annihilation conjecture (\lean{BernoulliSkeleton.lean}):
    that $\sum_{k \leq N} \mu(k) \cdot G_{\mathrm{arch}}(j,k) \to 0$
    sufficiently fast.
  \item PNT-level estimates on partial sums of
    $\sum \mu(k) f(k) / k$.
\end{enumerate}
These are deep number-theoretic inputs, not formalization gaps.
The Arakelov bridge identifies precisely \emph{where} the
mathematics must advance and \emph{why}.
